\documentclass[
	a4paper,
	landscape,
	%twoside,
	10pt,
	article
]{article}
\usepackage[
	a4paper,
	landscape,
	twocolumn,
	left=0.8cm,
	right=0.3cm,
% Weird top and bottom margins because of fancyhdr package
	top=1.8cm,
	bottom=-0.3cm,
	columnsep=1cm,
% Set margins on even and odd pages equal
	hmarginratio=1:1,
	asymmetric
]{geometry}
\usepackage[english]{babel}
\usepackage[utf8]{inputenc}
\usepackage{textcomp}
\usepackage{amsmath}
\usepackage{amsfonts}
\usepackage{graphicx}
\usepackage{float}
\usepackage{listings}
\usepackage{color}
\usepackage[colorinlistoftodos]{todonotes}
\usepackage[compact]{titlesec}		% shrink section whitespace
\usepackage{ifthen}
\usepackage{nicefrac}
\usepackage{hyperref}

%*************** Layout ***************
\setlength{\columnseprule}{0.2pt}
\newcommand{\latexcolumnseprulecolor}{\color{red}}
\titlespacing{\section}{0pt}{0pt}{0pt}
\sloppy

% A red divider in the middle of the page
\usepackage{etoolbox}
\makeatletter
\patchcmd\@outputdblcol{% find
	\normalcolor\vrule
}{% and replace by
	\latexcolumnseprulecolor\vrule
}{% success
}{% failure
	\@latex@warning{Patching \string\@outputdblcol\space failed}%
}
\makeatother


%*************** Title ***************

% all the \vspace are for reducing the vertical spacing
\title{
	\vspace{-5em}
	\includegraphics[scale=0.7]{./logo.eps}\\
	\vspace{-1.5em}
	Team Code Reference
	\vspace{-0.7em}
}
%\author{ Timon Knigge, Ragnar Groot Koerkamp, \& Harry Smit }
\author{
	\Large \textbf{ETH Team}\\
%	Timon Knigge, Ragnar Groot Koerkamp, \& Harry Smit
}
\date{
	\vspace{-0.7em}
	SWERC 2018
	\vspace{-1.9em}
}


%*************** Table of Contents ***************
\usepackage[toc]{multitoc}			% multicolumn toc
\usepackage{tocloft}				% to reduce toc spacing
\renewcommand*{\multicolumntoc}{2}
% reduce section spacing in toc
\setlength{\cftbeforesecskip}{-1pt}
\setlength{\cftbeforesubsecskip}{-1.5pt}
% remove the toc title
\makeatletter
\renewcommand{\@cftmaketoctitle}{}
\makeatother


%*************** Headings ***************
\usepackage{fancyhdr}
\pagestyle{fancy}
\fancyhead{}
\fancyfoot{}
\setlength{\headsep}{0.4em}
\setlength{\footskip}{0em}

% two sided
%\fancyhead[RE]{\bfseries Curiously Recurring \hspace{6.5em}}
%\fancyhead[LO]{\hspace{5em} Utrecht University}
%\fancyhead[LE,RO]{\thepage}
%\fancyhead[C]{\leftmark}

% one sided
\fancyhead[L]{\hspace{5em} ETH Z\"urich \phantom{-} \bfseries
ETH Team}
\fancyhead[R]{\thepage \hspace{0.5em}}
\fancyhead[C]{\leftmark}

%*************** Code highlighting ***************
\lstset{
	backgroundcolor=\color{white},
	tabsize=4,
	language=C++,
	basicstyle=\footnotesize\ttfamily,
	frame=lines,
	numbers=left,
	numberstyle=\tiny,
	numbersep=5pt,
	breaklines=true,
	keywordstyle=\color[rgb]{0, 0, 1},
	commentstyle=\color[rgb]{0, 0.5, 0},
	stringstyle=\color{red}
}


%*************** Section entries ***************
% \entry{name}{description}{snippet location}{complexity}{dependencies}
\newcommand{\entry}[3]{
	\subsection{#1}
	#2
	\ifthenelse{\equal{#3}{}}{}{\lstinputlisting[firstline=2]{#3}}
}
\newcommand{\otherentry}[3]{
	\subsection{#1}
	#2
	\lstinputlisting[language=]{#3}
}


%*************** Begin document ***************
\begin{document}


%*************** Reduce align spacing ***************
\setlength{\abovedisplayskip}{0pt}
\setlength{\belowdisplayskip}{0pt}
\setlength{\abovedisplayshortskip}{0pt}
\setlength{\belowdisplayshortskip}{0pt}

%*************** Titlepage ***************
{\let\newpage\relax\maketitle}
\tableofcontents
\thispagestyle{empty}
\newpage

%*************** Contents ***************

\section{Datastructures}

\entry{Union Find}
{}
{./snippets/datastructures/unionfind.cpp}

\entry{Fenwick Tree}
{Can be generalized to arbitrary dimensions by duplicating loops.}
{./snippets/datastructures/fenwick.cpp}

\entry{Segment Tree}{Works bottom up, so not suitable for lazy propagation.}
{./snippets/datastructures/segmenttree.cpp}

\entry{Lazy Dynamic Segment Tree}
{Here \texttt{T} is the data type and \texttt{U} the lazy update type.
Implement all helper functions.}
{./snippets/datastructures/segmenttree_dynamic.cpp}

\entry{Heavy-Light decomposition}{}
{./snippets/datastructures/heavylight.cpp}


\entry{Sequence}{This is essentially a treap, but it needs special
	functionality for the Euler Tour Tree}
{./snippets/datastructures/sequence.cpp}

\entry{Euler Tour Tree}{Maintain information about the trees using
	the underlying \texttt{seq} datastructure.}
{./snippets/datastructures/euler-tour-tree.cpp}

\entry{Suffix Array}{An $O(n \log n)$ implementation. Note that the resulting
	array maps values to their position in the suffix array, so invert if
	necessary.}
{./snippets/datastructures/suffixarray.cpp}

\entry{Convex Hull Set}{}{./snippets/datastructures/convex-hull-set.cpp}

\entry{GNU Built-in datastructures}
{These require \texttt{gnu} so check in test session.}
{./snippets/datastructures/builtin.cpp}

\section{Combinatorics and Exact Algorithms}

\subsection{Formulae}
\subsubsection*{Permutations}
The number of derangements $D_n$, $n$-permutations without a fixed point,
satisfies $D_{n+1} = n(D_n + D_{n-1}) = (n+1)D_n + (-1)^{n+1}$. If
$g_S(n)$ counts the number of $n$-permutations whose cycle lengths are in
$S$, then
		$$\sum_{n=0}^\infty g_S(n)\frac{x^n}{n!} = \exp\Big(\sum_{k\in S}
			\frac{x^k}{k}\Big)$$

Stirling numbers of the first kind $S_1(n, k)$ count permutations on $n$ items
with $k$ cycles. $S_1(n, k) = S_1(n-1, k-1) + (n-1)S_1(n-1, k)$ with
$S_1(0, 0) = 1$. Note $\sum_{k=0}^n S_1(n, k)x^k = x(x+1)\dots(x+n-1)$.

Eulerian numbers $E(n, k)$ count the number of permutations on $n$ elments,
with exactly $k$ elements that are greater than their predecessor, i.e.
$k$ `ascents'. We have $E(n, k) = (n-k)E(n-1,k-1)+(k+1)E(n-1,k)$ with
$E(n, 0) = E(n, n-1)=1$.

\subsubsection*{Binomials and other partitionings}
We have $\binom{n}{k} = \binom{n-1}{k}+\binom{n-1}{k-1} =
	\prod_{i=1}^k \frac{n-i+1}{i}$. This last product may be computed
incrementally since any product of $k'$ consecutive values is divisibleby
$k'!$.

Hockeystick identity: $\sum_{k=r}^n \binom{k}{r} = \binom{n+1}{r+1}$.

For $n, m \geq 0$ and $p$ prime. Write $n, m$ in base $p$, i.e.
$n = n_k p^k + \dots + n_1 p + n_0$ and $m = m_k p^k + \dots m_1 p + m_0$. Then
by Lucas theorem we have $\binom{n}{m} \equiv \prod_{i=0}^k \binom{n_i}{m_i}
	\mod p$, with the convention that $n_i < m_i \implies \binom{n_i}{m_i} =0$.

We can put $n$ indistinguishable items into $k$ distinguishable bins in exactly
$\binom{n+k-1}{k-1}$ ways, or if we require the bins be non-empty, in
$\binom{n-1}{k-1}$ ways.

Stirling numbers of the second kind $S_2(n, k)$ count partitions of $n$
distinct elements into exactly $k$ non-empty groups.
$S_2(n, k) = S_2(n-1, k-1) + kS_2(n-1, k)$ with $S_2(n, 1) = S_2(n, n) = 1$ and
$$S_2(n, k) = \frac{1}{k!}\sum_{i=0}^k (-1)^{k-i}\binom{k}{i}i^n$$

Bell numbers $B_n$ count arbitrary partitions of $n$ distinct elements, i.e.
$B_n = \sum_k S_2(n, k)$.

\subsubsection*{Catalan numbers}
Catalan numbers $C_n$ satisfy $C_n = \frac{1}{2n+1}\binom{2n}{n}$ as well
as $C_0 = 1, C_{n+1} = \sum_i C_i C_{n-i}$. These count, among other things:
\begin{itemize}
	\item Valid parenthesis sequences of length $2n$.
	\item Sub-diagonal monotone paths from $(0, 0)$ to $(n, n)$ in the
		$n\times n$ grid.
	\item Complete binary trees (i.e. no vertices with exactly $1$ child)
		with $n+1$ leaves.
	\item Triangulations of a convex polygon with $n+2$ sides.
	\item Stack-sortable permutations of size $n$.
\end{itemize}

\subsubsection*{Trees}
There are $n^{n-2}$ labeled unrooted (cq. $n^{n-1}$ rooted) trees on $n$
vertices. A bijection between a tree and it's Pr\"ufer sequence
$(a_1, a_2, \dots, a_{n-2})$is given by:
\begin{enumerate}
	\item[$\leftarrow$] Given the Pr\"ufer sequence, let $d_i$ be $1$ plus the
		number of occurences of $i$ in $a$. Now for each $a_j$ for $j=1\dots$,
		find the lowest
		numbered node $k$ with degree $1$ and add $\{a_j, k\}$, and decrement
		the degrees $d_{a_j}, d_k$. At the end two degree $1$ nodes remain,
		connect them.
	\item[$\rightarrow$] Iterate, at step $i=1\dots$ pick the leaf with the
		smallest label and set the $i^{th}$ element of $a$ to be its
		neighbour.
\end{enumerate}

\subsection{Burnside's lemma}


\subsection{Inclusion-Exclusion}

For appropriate $f$ compute $\sum_{S\subseteq T} (-1)^{|T\setminus S} f(S)$.
In some contexts we might use Stirling numbers, not binomial coefficients!

Some useful applications:
\begin{enumerate}
	\item[Graph coloring] Let $I(S)$ count the number of independent sets
		contained in $S \subseteq V$ ($I(\emptyset) = 1$,
		$I(S) = I(S\setminus v) + I(S\setminus N(v))$). Let
		$c_k = \sum_{S\subseteq V} (-1)^{|V\setminus S|} I(S)$. Then $V$
		is $k$-colorable iff $v > 0$.
	\item[TODO]
\end{enumerate}

\subsection{Stable Marriage Algorithm}
Given $n$ men and $n$ women, with each man ranking each woman by preference
(and vice versa), we can find a stable marriage by trying to match each man
to his most preferred woman that wasn't tried before. Runs in $O(n^2)$.

\section{Dynamic Programming}

\entry{LIS}
{For the actual sequence, store parent pointers. Note that longest common
subsequence can be reduced to LIS with unique elements.}
{./snippets/dp/lis.cpp}

\entry{All Nearest Smaller Values}{}{./snippets/dp/ansv.cpp}

\subsection{DP Optimalizations}
\begin{enumerate}
	\item[CH] Recurrences are generally of the form
		$dp_i = dp_j \cdot a_i + b_j$ or equivalent. Insert all lines as
		functions of $a_i$ into the convex hull set (see datastructures) to
		query in logarithmic time.
	\item[DC] Recurrences are generally of the form
		$dp_{k,i} = \max_{j<i} f(dp_{k-1,j}, i, j, k)$. If the $\arg\max$ is
		monotonic in $i$, do divide and conquer by first finding the answer
		for $i = \frac{n}{2}$, and then recursing for larger and smaller $i$.
	\item[Knuth] Recurrences are generally of form similar to
		$$dp_{i,j} = min_{i\leq k<j} \big(f(i, j) + dp_{i,k} + dp_{k+1,j}\big)$$
		Let $a_{i,j}$ be the $\text{argmin}$ for $dp_{i, j}$, then if
		$a_{i,j-1} \leq a_{i, j} \leq a_{i+1,j}$, for a fixed $len > 0$ the
		ranges we have to scan for all $dp_{i, i+len}$ are (almost) disjoint,
		and we can solve the recurrence in $O(n^2)$ time.

		A sufficient condition for applicability on $f$ is, for
		$i\leq j\leq k\leq l$, monotonicity: $f(j,k) \leq f(i, l)$, and
		the quadrangle inequality: $f(i,k)+f(j,l)\leq f(i,l)+f(j,k)$.
\end{enumerate}

\section{Geometry}

\subsection{Formulae}
Line equations $ax + by = e$. For a line going through points $(x_1, y_1)$ and
$(x_2, y_2)$, note that $n = (y_1 - y_2, x_2 - x_1)$ is a normal vector, so
the line equation is given by $n \cdot p = n \cdot (x_1, y_1)$ for
$p \in.\mathbb{R}^2$.
$$\begin{aligned}ax+by=e\\cx+dy=f\end{aligned}
\,\to\,
\begin{aligned}x=\dfrac{ed-bf}{ad-bc}\\y=\dfrac{af-ec}{ad-bc}\end{aligned}$$

For a general equation of the form $Ax = b$, we have $x_i = \det A_i / \det A$ with $A_i$ equal to $A$ with the $i^{th}$ column replaced by $b$.
\end{document}
